% Example generations on the two reasoning tasks (verbatim narrator outputs, greedy).
% Galaxy: galaxy_id 601201, narrator trained with LoRA and no head token
%   (MultiLLM cache/galaxy_gen_eval/generations/lora_nohead_xmodal_eval.jsonl).
% Stars: test group population_451_N4 (MultiLLM cache/stellar_gen_eval_full/
%   generations/{astrosage8b,flow_astrosage8b_template}.jsonl).
\newcommand{\gOK}{{\color{teal!80!black}\,$\checkmark$}}
\newcommand{\gNO}{{\color{coral}\,$\times$}}
\newcommand{\hGood}[1]{{\setlength{\fboxsep}{0.6pt}\colorbox{paleteal}{\strut #1}}}
\newcommand{\hBad}[1]{{\setlength{\fboxsep}{0.6pt}\colorbox{palecoral}{\strut #1}}}
\tcbset{exbox/.style={colback=white, boxrule=0.6pt, arc=1.5pt,
  left=3pt, right=3pt, top=2pt, bottom=2pt, fontupper=\footnotesize,
  fonttitle=\footnotesize\bfseries, coltitle=deepteal,
  toptitle=0.6pt, bottomtitle=0.6pt}}
\begin{minipage}[t]{0.47\textwidth}\vspace{0pt}
  \begin{tcolorbox}[exbox, colframe=neutral, colbacktitle=palegrey,
      title={Galaxy: does its shape agree with its history?}]
    \raggedright
    \textbf{Catalog.} Late-type disk ($P{=}0.75$); quiescent
    ($\log \mathrm{sSFR} = -11.6$); declining star-formation history.
    A disk that has stopped forming stars: shape and history
    \emph{disagree}.
  \end{tcolorbox}\vspace{-2pt}
  \begin{tcolorbox}[exbox, colframe=teal!80!black, colbacktitle=paleteal,
      title={Narrator (no head token)\gOK}]
    \raggedright
    A \hGood{quiescent} \hGood{late-type disk} --- a
    \hGood{cross-modal outlier}, because disks rarely quench while retaining
    their morphology. Its star-formation history is \hGood{declining}.
  \end{tcolorbox}
  {\scriptsize\color{neutral} Shape, star formation and verdict all match the
  catalog.}
\end{minipage}\hfill
\begin{minipage}[t]{0.5\textwidth}\vspace{0pt}
  \begin{tcolorbox}[exbox, colframe=neutral, colbacktitle=palegrey,
      title={Stars: which one does not belong?}]
    \centering
    \begin{tabular}{@{}clccc@{}}
      star & class & $T_\mathrm{eff}$ [K] & $\log g$ & $R$ [$R_\odot$] \\
      \hline
      0 & dwarf    & 5778 & 4.39 & 1.05 \\
      \hBad{1} & \hBad{\textbf{giant}} & \hBad{4662} & \hBad{\textbf{2.45}} & \hBad{\textbf{11.14}} \\
      2 & subgiant & 5628 & 3.89 & 2.00 \\
      3 & dwarf    & 5394 & 4.55 & 0.81 \\
    \end{tabular}
  \end{tcolorbox}\vspace{-2pt}
  \begin{tcolorbox}[exbox, colframe=coral, colbacktitle=palecoral,
      title={Narrator, no head\gNO}]
    \raggedright
    These 4 stars form a homogeneous population of field stars.
    \hBad{No anomalies detected} along the field-star axis.
  \end{tcolorbox}\vspace{-2pt}
  \begin{tcolorbox}[exbox, colframe=teal!80!black, colbacktitle=paleteal,
      title={Narrator $+$ head\gOK}]
    \raggedright
    Sample N=4: single stars. \hGood{index 1: a giant} [logg = 2.44 dex,
    R = 11.14 Rsun, logL = 1.72 log(Lsun), giant].
  \end{tcolorbox}
\end{minipage}
